%! Dimensions of the Four Subspaces — Unit 3, Lesson 3.5 — Group Activity
\documentclass[10pt]{article}
\usepackage{linalg-article}
\usepackage{linalg-boxes}
\newcommand{\TallMath}[1]{$\displaystyle #1\rule[-1.4em]{0pt}{3.2em}$}
\newcommand{\vv}{\mathbf{v}}
\newcommand{\ww}{\mathbf{w}}
\newcommand{\xx}{\mathbf{x}}
\newcommand{\yy}{\mathbf{y}}
\newcommand{\svec}{\mathbf{s}}
\newcommand{\zero}{\mathbf{0}}
\newcommand{\T}{\mathsf{T}}

\begin{document}

\pageheader{Unit 3, Lesson 3.5}{Group Activity}
\namepartnerperiod

\vspace{0.10in}

\begin{headlinebox}{blush}
\small\textbf{The Four Subspaces.} An $m\times n$ matrix $A$ of rank $r$ has four fundamental subspaces. In
the \emph{input room} $\mathbb{R}^n$: the row space $C(A^{\T})$ ($\dim r$) and the nullspace $N(A)$
($\dim n-r$). In the \emph{output room} $\mathbb{R}^m$: the column space $C(A)$ ($\dim r$) and the left
nullspace $N(A^{\T})$ ($\dim m-r$). \emph{One} reduction of $A$ gives a basis for all four. Work with your
partner and \emph{justify} every answer.
\end{headlinebox}

\vspace{0.10in}

\begin{tcolorbox}[colback=white, colframe=black!40, title=\textbf{Tier R --- Remediate},
    fonttitle=\bfseries, arc=1.5mm, left=3mm, right=3mm, top=2mm, bottom=2mm, breakable]
The matrix below is already reduced: $R=\begin{bmatrix} 1 & 0 & 2 \\ 0 & 1 & 4 \\ 0 & 0 & 0 \end{bmatrix}$.
\begin{enumerate}[leftmargin=1.7em, itemsep=8pt]
  \item \textbf{Three numbers.} Rows $m=\blank{1.0cm}$\quad columns $n=\blank{1.0cm}$\quad pivots $r=\blank{1.0cm}$
  \item \textbf{Which room?} Circle the room each subspace lives in.
    \begin{enumerate}[label=(\alph*), leftmargin=1.8em, itemsep=4pt]
      \item Column space $C(A)$: \quad $\mathbb{R}^m$ / $\mathbb{R}^n$ \qquad\qquad
            (c)\; Nullspace $N(A)$: \quad $\mathbb{R}^m$ / $\mathbb{R}^n$
      \item Row space $C(A^{\T})$: \quad $\mathbb{R}^m$ / $\mathbb{R}^n$ \qquad\quad
            (d)\; Left nullspace $N(A^{\T})$: \quad $\mathbb{R}^m$ / $\mathbb{R}^n$
    \end{enumerate}
  \item \textbf{Four dimensions from $r$.}
        $\dim C(A)=\blank{0.9cm}$\quad $\dim C(A^{\T})=\blank{0.9cm}$\quad
        $\dim N(A)=\blank{0.9cm}$\quad $\dim N(A^{\T})=\blank{0.9cm}$
\end{enumerate}
\end{tcolorbox}

\begin{tcolorbox}[colback=white, colframe=black!40, title=\textbf{Tier A --- Approaching Proficiency},
    fonttitle=\bfseries, arc=1.5mm, left=3mm, right=3mm, top=2mm, bottom=2mm, breakable]
For each matrix: reduce, then give a basis \emph{and} dimension for all four subspaces, and check both
counting rules $r+(n-r)=n$ and $r+(m-r)=m$. Note that the two rooms $\mathbb{R}^m$ and $\mathbb{R}^n$ have
\emph{different} sizes here.
\begin{enumerate}[leftmargin=1.7em, itemsep=9pt]
  \item $A=\begin{bmatrix} 1 & 1 & 2 \\ 1 & 2 & 3 \end{bmatrix}$ \quad (\emph{$2\times3$; watch what the left nullspace does}). \writelines{4}
  \item $A=\begin{bmatrix} 1 & 1 & 2 & 1 \\ 1 & 2 & 3 & 1 \\ 2 & 3 & 5 & 2 \end{bmatrix}$ \quad (\emph{$3\times4$; row~3 $=$ row~1 $+$ row~2}). \writelines{5}
\end{enumerate}
\end{tcolorbox}

\begin{tcolorbox}[colback=white, colframe=black!40, title=\textbf{Tier E --- Extension},
    fonttitle=\bfseries, arc=1.5mm, left=3mm, right=3mm, top=2mm, bottom=2mm, breakable]
\textbf{Why the counting works.}
\begin{enumerate}[leftmargin=1.7em, itemsep=9pt]
  \item \textbf{Row rank $=$ column rank.} A reduction of $A$ has exactly $r$ pivots. Explain why that one
        number counts \emph{both} the independent columns (the pivot columns) and the independent rows (the
        nonzero rows of the reduced form) --- so $\dim C(A)=\dim C(A^{\T})=r$. \writelines{2}
  \item \textbf{When is the left nullspace just $\{\zero\}$?} Explain why $N(A^{\T})=\{\zero\}$ exactly when
        $r=m$ (``full row rank''), and what that says about redundancy among the rows. \writelines{2}
  \item \textbf{Read the row dependency.} For $A=\begin{bmatrix} 1 & 1 & 2 \\ 2 & 1 & 3 \\ 3 & 2 & 5 \end{bmatrix}$
        a left-nullspace vector is $\yy=\begin{bsmallmatrix} 1\\1\\-1 \end{bsmallmatrix}$. Write the exact
        combination of the rows of $A$ that this says equals the zero row. \writelines{2}
\end{enumerate}
\end{tcolorbox}

\end{document}
